\section{Proof details and complete ledgers}
\label{app:proofs}

\subsection{Bounds and translated witnesses}

Right multiplication in $P_r$ is injective, so $S$ and $T$ are isometries.
This gives $\|\Aplus\|\le a+b$ and
$\|\Asym\|\le2(a+b)$.  The symmetrized graph has degree at most four.  From
an oriented edge there are at most three nonbacktracking successors, each of
weight at most $\max(a,b)$; the Schur test gives
$\|B\|\le3\max(a,b)$.

For $x_j=(0,j)$, the supports of $\Aplus\delta_{x_j}$ are pairwise disjoint.
For $\Asym$, the subsequence $x_{4j}$ separates predecessor and successor
supports.  For $B$, use the oriented $V$ edges at levels $j$; the $U$ and $V$
continuations are distinct and remain disjoint across translated levels.
These are the orthogonal witnesses used in
\cref{prop:noncompactness}.

\subsection{Primitivity and infinite multiplicity}

The cyclic word $v u\bar v\bar u^r$ has one positive $v$ and one negative
$v$.  It is freely and cyclically reduced: no adjacent pair, including the
cyclic join, consists of an edge and its reverse.  If it were a $k$-fold
cyclic power for $k\ge2$, the count of positive $v$ occurrences would be
divisible by $k$, a contradiction.  Choosing base vertices at sufficiently
separated height levels yields infinitely many edge-disjoint translates.

Ordinary closed words at a base $x$ are exactly admissible identity words in
$G_r=\Z[1/r]\rtimes\Z$.  This proves completeness of the word-level
classification: free and cyclic reduction impose the Hashimoto rule, and
the usual least-period test imposes temporal primitivity.  Infinite
translate multiplicity is part of the obstruction.

\subsection{Matrix local-factor boundary}

For a $d$-dimensional holonomy $M$ with eigenvalues $\lambda_j$,
\[
\det(I-tM)=\prod_{j=1}^d(1-t\lambda_j).
\]
This polynomial is identically one precisely when every eigenvalue is zero,
equivalently when $M$ is nilpotent.  An invertible or unitary holonomy cannot
delete its complete Euler factor.  Cancellation between different primitive
orbits is a different coefficientwise statement and is not supplied by this
lemma.
